\subsection{Hard-family local facts and class membership}
\label{app:lower-proof-family}

This subsection proves the one-edge calculations in
\cref{app:lem:hard-local-facts}: positivity and common marginals, the edge
response, the selected-oracle error, and the one-edge Hellinger loss.  The
same calculations then verify \cref{app:prop:hard-family-inclusion}, placing
the finite hard family inside the shared target/oracle class used by both
bounds.

\paragraph{Intuition.}
There are two different error scales.  Removing one dependence edge costs
$\HardEdgeDistance\simeq\EdgeCoupling^2$, whereas replacing a nontrigger
conditional by the neutral marginal costs at most order
$\EdgeCoupling^2\TriggerMass$.  The extra factor $\TriggerMass$ lets the
frozen oracle hide failed trigger tests while preserving a larger loss
when two unresolved endpoints are committed together.

For the moment calculations, if
\(A\sim\HardMarginal\), then the trigger indicator at position \(i\) is
Bernoulli with mean \(\TriggerMass\).  Hence
\begin{equation}
 \mathbb E\TriggerFeature{i}(A)=0,
 \qquad
 \mathbb E\TriggerFeature{i}(A)^2=\TriggerVariance,
 \qquad
 \mathbb E|\TriggerFeature{i}(A)|=2\TriggerVariance.
 \label{app:eq:hard-feature-moments}
\end{equation}
Indeed, the feature is $1-\TriggerMass$ with probability $\TriggerMass$
and $-\TriggerMass$ otherwise, so
\[
\begin{aligned}
 \mathbb E\TriggerFeature{i}
 &=\TriggerMass(1-\TriggerMass)
   -(1-\TriggerMass)\TriggerMass=0,\\
 \mathbb E\TriggerFeature{i}^2
 &=\TriggerMass(1-\TriggerMass)^2
   +(1-\TriggerMass)\TriggerMass^2
   =\TriggerMass(1-\TriggerMass),\\
 \mathbb E|\TriggerFeature{i}|
 &=\TriggerMass(1-\TriggerMass)
   +(1-\TriggerMass)\TriggerMass
   =2\TriggerMass(1-\TriggerMass).
\end{aligned}
\]

\begin{lemma}[Local facts for the hard matching family]
\label{app:lem:hard-local-facts}
Under \eqref{app:eq:hard-basic-conditions}, every member of
\(\HardFamily[\HardVocabSize,\HardBankSize,
\TriggerMass,\EdgeCoupling]\) has the following properties.
\begin{enumerate}
 \item \textup{\textbf{Positivity, marginals, and factorization.}}
 Each edge law in \eqref{app:eq:hard-edge-law} is strictly positive,
 has both marginals equal to \(\HardMarginal\), and the product law satisfies
 the structural forest assumptions.
 \item \textup{\textbf{Edge response.}}
 For either orientation \(i\to j\) of a matching edge, the largest TV
 change in the endpoint kernel equals \(\EdgeCoupling\).
 \item \textup{\textbf{Oracle error.}}
 Let $j$ be any readout, let $i:=\MateInMatching{j}$, and let
 $a\in\Vocab\setminus\{\HardCandidate{\TriggerIndex{i}}\}$ be a revealed
 nontrigger value at $i$.  The selected oracle is exact outside this case,
 while in this case
 \begin{equation}
  \sqHellinger\bigl(
   \HardKernel{i}{j}{I}(\cdot\mid a),\HardMarginal
  \bigr)
  \le
  \frac{\EdgeCoupling^2\TriggerMass}
       {2(1-\TriggerMass)}
  \le \EdgeCoupling^2\TriggerMass.
  \label{app:eq:hard-local-oracle-error}
 \end{equation}
 \item \textup{\textbf{One-edge Hellinger loss.}}
 For every matching edge \(e\), uniformly over its two private trigger
 indices, the value below depends only on
 \((\TriggerMass,\EdgeCoupling)\), not on \(e\) or \(I\):
 \begin{equation}
  \frac{\EdgeCoupling^2}{12}
  \le
  \HardEdgeDistance
  :=\sqHellinger\bigl(
     \HardEdgeLaw{e}{I},\HardIndependentEdgeLaw{e}
    \bigr)
  \le\frac{\EdgeCoupling^2}{2}.
  \label{app:eq:hard-edge-distance}
 \end{equation}
\end{enumerate}
\end{lemma}

\begin{proof}

\emph{1. Positivity, common marginals, and forest factorization.}
Fix \(e=\{i,j\}\).  Summing \eqref{app:eq:hard-edge-law} over its second
coordinate and using \eqref{app:eq:hard-feature-moments} gives
\[
 \sum_b\HardEdgeLaw{e}{I}(a,b)
 =\HardMarginal(a)
  \left[1+
   \EdgeCoupling\frac{\TriggerFeature{i}(a)}{\TriggerVariance}
   \mathbb E\TriggerFeature{j}(B)
  \right]
 =\HardMarginal(a),
\]
where \(B\sim\HardMarginal\).  The other marginal follows symmetrically.
The likelihood ratio with respect to
\(\HardIndependentEdgeLaw{e}\) takes the three values
\begin{equation}
 1+\EdgeCoupling
   \frac{\TriggerFeature{i}(a)\TriggerFeature{j}(b)}
        {\TriggerVariance}
 =
 \begin{cases}
  1+\EdgeCoupling(1-\TriggerMass)/\TriggerMass,
    &a,b\text{ are their respective triggers},\\
  1-\EdgeCoupling,
    &\text{exactly one of }a,b\text{ is its trigger},\\
  1+\EdgeCoupling\TriggerMass/(1-\TriggerMass),
    &a,b\text{ are both nontriggers}.
 \end{cases}
 \label{app:eq:hard-likelihood-values}
\end{equation}
The conditions
\(0<\EdgeCoupling\le\TriggerMass\le1/2\) place every value in
\([1/2,2]\), proving strict positivity.  The marginal identity also gives
the endpoint compatibility
\[
 \HardMarginal(a)\HardKernel{i}{j}{I}(b\mid a)
 =\HardEdgeLaw{e}{I}(a,b)
 =\HardMarginal(b)\HardKernel{j}{i}{I}(a\mid b).
\]
Every component is one edge, so multiplying these identities over the hidden
matching proves the required forest factorization for either root choice.


\emph{2. Directed endpoint response.}
For source values \(a,a'\),
\begin{align}
 &\dTV\bigl(
   \HardKernel{i}{j}{I}(\cdot\mid a),
   \HardKernel{i}{j}{I}(\cdot\mid a')
  \bigr)\notag\\
 &\qquad=
 \frac{\EdgeCoupling
       |\TriggerFeature{i}(a)-\TriggerFeature{i}(a')|}
      {2\TriggerVariance}
 \sum_b\HardMarginal(b)|\TriggerFeature{j}(b)|
 =\EdgeCoupling
  |\TriggerFeature{i}(a)-\TriggerFeature{i}(a')|.
 \label{app:eq:hard-directed-response}
\end{align}
The feature equals \(1-\TriggerMass\) at the unique trigger and
\(-\TriggerMass\) elsewhere.  Its largest difference is therefore one, so
the maximum in \eqref{app:eq:hard-directed-response} is exactly
\(\EdgeCoupling\).


\emph{3. Oracle error at a revealed nontrigger.}
Because the target is a product over matching edges, the exact row at a
readout \(j\) is \(\HardMarginal\) when its mate is masked and is
\(\HardKernel{i}{j}{I}(\cdot\mid y_i)\) when the mate \(i\) is revealed.
Thus \eqref{app:eq:hard-oracle} is exact in its first two cases.  In the third
case, $\TriggerFeature{i}(y_i)=-\TriggerMass$, so the likelihood ratio is
\[
 \frac{\HardKernel{i}{j}{I}(b\mid y_i)}{\HardMarginal(b)}
 =1-\frac{\EdgeCoupling}{1-\TriggerMass}\TriggerFeature{j}(b).
\]
Substitution into the definition of chi-square divergence gives
\begin{align}
 \ChiSquareDivergence{
  \HardKernel{i}{j}{I}(\cdot\mid y_i)}{\HardMarginal}
 &=\sum_b\HardMarginal(b)
       \left(\frac{\EdgeCoupling\TriggerFeature{j}(b)}
                   {1-\TriggerMass}\right)^2\notag\\
 &=\frac{\EdgeCoupling^2}{(1-\TriggerMass)^2}
       \TriggerMass(1-\TriggerMass)
 =\EdgeCoupling^2\frac{\TriggerMass}{1-\TriggerMass}.
 \label{app:eq:hard-oracle-chi-square}
\end{align}
For a likelihood ratio \(r\),
\((\sqrt r-1)^2\le(r-1)^2\).  The normalization in
\eqref{app:eq:distances} therefore implies
\(\sqHellinger(p,q)\le\ChiSquareDivergence{p}{q}/2\).  Combining this with
\eqref{app:eq:hard-oracle-chi-square} and
\(\TriggerMass\le1/2\) proves \eqref{app:eq:hard-local-oracle-error}.


\emph{4. One-edge squared-Hellinger loss.}
Finally, put
\[
 Z_e(a,b):=
 \frac{\TriggerFeature{i}(a)\TriggerFeature{j}(b)}
      {\TriggerVariance}.
\]
Under \(\HardIndependentEdgeLaw{e}\), each endpoint independently equals
its trigger token with probability \(\TriggerMass\).  Thus
\[
 Z_e=
 \begin{cases}
 (1-\TriggerMass)/\TriggerMass,
   &\text{with probability }\TriggerMass^2
      \quad\text{(both triggers)},\\
 -1,
   &\text{with probability }2\TriggerMass(1-\TriggerMass)
      \quad\text{(one trigger)},\\
 \TriggerMass/(1-\TriggerMass),
   &\text{with probability }(1-\TriggerMass)^2
      \quad\text{(neither trigger)}.
 \end{cases}
 \]
This law depends only on \(\TriggerMass\), not on the edge or its trigger
indices.
\Eqref{app:eq:hard-feature-moments} and independence under
\(\HardIndependentEdgeLaw{e}\) give
\[
 \mathbb E_{\HardIndependentEdgeLaw{e}}Z_e
 =\frac{\mathbb E\TriggerFeature{i}\,
        \mathbb E\TriggerFeature{j}}{\TriggerVariance}=0,\qquad
 \mathbb E_{\HardIndependentEdgeLaw{e}}Z_e^2
 =\frac{\mathbb E\TriggerFeature{i}^2\,
        \mathbb E\TriggerFeature{j}^2}{(\TriggerVariance)^2}=1.
\]  Since
\(\mathrm d\HardEdgeLaw{e}{I}/\mathrm d\HardIndependentEdgeLaw{e}
=1+\EdgeCoupling Z_e\),
the definition of squared Hellinger distance gives the first line below.
For the second, multiply
$\sqrt{1+\EdgeCoupling Z_e}-1$ by
$(\sqrt{1+\EdgeCoupling Z_e}+1)/
 (\sqrt{1+\EdgeCoupling Z_e}+1)$:
\begin{align}
 \HardEdgeDistance
 &=\frac12\mathbb E_{\HardIndependentEdgeLaw{e}}
       \bigl(\sqrt{1+\EdgeCoupling Z_e}-1\bigr)^2\notag\\
 &=\frac{\EdgeCoupling^2}{2}
  \mathbb E_{\HardIndependentEdgeLaw{e}}
  \frac{Z_e^2}
       {(\sqrt{1+\EdgeCoupling Z_e}+1)^2}.
 \label{app:eq:hard-edge-distance-formula}
\end{align}
Consequently \(\HardEdgeDistance\) depends only on
\((\TriggerMass,\EdgeCoupling)\), so it is the same for every matching
component and every choice of its trigger indices.
By \eqref{app:eq:hard-likelihood-values}, the denominator lies between
one and \((\sqrt2+1)^2<6\).
\Eqref{app:eq:hard-edge-distance} follows.
\end{proof}

\paragraph{Intuition for class inclusion.}
The four checks have separate jobs:
$\HardVocabSize\TriggerMass>1$ puts the candidate bank above the background,
$\TriggerMass\le\HardBankSize^{-\RankTailExponent}$ gives the frequency
envelope, $\EdgeSignal\le\EdgeCoupling$ gives the edge signal, and
$\EdgeCoupling^2\TriggerMass\le(\OracleRadius)^2/(2N)$ hides nontrigger replies
within the oracle budget.  In the main calibration,
$\EdgeCoupling=\TriggerMass=\EdgeSignal$, so the last check becomes the
cubic constraint $\EdgeSignal^3\le(\OracleRadius)^2/(2N)$.

\begin{proposition}[Finite hard-family inclusion]
\label{app:prop:hard-family-inclusion}
In addition to \eqref{app:eq:hard-basic-conditions}, let
$\RankTailExponent>1$, use the same public vocabulary of size $\VocabSize$,
set
$\ResponseExponent=\ResponseConstant=\FrequencyConstant=1$, and suppose
\begin{equation}
 \HardVocabSize\TriggerMass>1,
 \qquad
 \TriggerMass\le\HardBankSize^{-\RankTailExponent},
 \qquad
 \EdgeSignal\le\EdgeCoupling,
 \qquad
 \EdgeCoupling^2\TriggerMass\le\frac{(\OracleRadius)^2}{2N}.
 \label{app:eq:finite-hard-inclusion}
\end{equation}
Then every target in
$\HardFamily[\HardVocabSize,\HardBankSize,\TriggerMass,\EdgeCoupling]$
satisfies the structural and response assumptions in
Appendix~\ref{app:setup}, as well as the finite rank--frequency inequality in
(RF\textcolor{red}{; \cref{app:eq:RF}}), and its selected oracle \eqref{app:eq:hard-oracle} belongs to
$\OracleClass(\HardTargetLaw{I};\OracleRadius)$.  If the public vocabulary
sequence also obeys the vocabulary-growth condition in (RF\textcolor{red}{; \cref{app:ass:rate-regime}}), then the full
(RF\textcolor{red}{; \cref{app:ass:rate-regime}}) assumption holds.  Under that additional asymptotic condition, this
matching family is a contained subfamily of the same public problem class
used by the upper theorem.
\end{proposition}

\begin{proof}[Proof of \cref{app:prop:hard-family-inclusion}]

\emph{1. Forest structure and the edge-signal condition (UEN\textcolor{red}{; \cref{app:eq:UEN}}).}
Apply \cref{app:lem:hard-local-facts}.  Its first conclusion proves
strict positivity and Assumption~\ref{app:ass:forest-structure}.  Complete
boundary assignments outside a matching edge do not alter that edge's law,
so \eqref{app:eq:hard-directed-response} makes both directed intrinsic
responses exactly \(\EdgeCoupling\).  The condition
\(\EdgeSignal\le\EdgeCoupling\) therefore proves (UEN\textcolor{red}{; \cref{app:eq:UEN}}).


\emph{2. Response-tail regularity (RT\textcolor{red}{; \cref{app:eq:RT}}).}
For (RT\textcolor{red}{; \cref{app:eq:RT}}), first consider distinct positions \(i,j\) that are not matched.
Conditioning on a complete boundary fixes the mates of \(i\) and \(j\), and
the product over matching components makes
\(\BoundaryRow{i}{j}{a}{z}\) independent of \(a\).  The response diameter is
therefore zero.

It remains to consider a matching edge.  Fix a threshold
\(t\).
\begin{itemize}
\item If \(t<\TriggerMass\), the marginal tail
contains no candidate token and hence no trigger.  Every source value in the
tail has the same centered feature \(-\TriggerMass\), so its response
diameter is zero.
\item If \(t\ge\TriggerMass\), the response diameter is at most
\(\EdgeCoupling\le\TriggerMass\le t\).
\end{itemize}
Thus (RT\textcolor{red}{; \cref{app:eq:RT}}) holds with
\(\ResponseExponent=\ResponseConstant=1\).


\emph{3. Rank--frequency regularity (RF\textcolor{red}{; \cref{app:eq:RF}}).}
The inequality \(\HardVocabSize\TriggerMass>1\) is equivalent to the
background mass in \eqref{app:eq:hard-common-marginal} being strictly less
than $\HardVocabSize^{-1}$, since
\[
 \frac{1-\HardBankSize\TriggerMass}{\HardVocabSize-\HardBankSize}
 <\frac1{\HardVocabSize}
 \quad\Longleftrightarrow\quad
 \HardVocabSize-\HardVocabSize\HardBankSize\TriggerMass
 <\HardVocabSize-\HardBankSize
 \quad\Longleftrightarrow\quad
 \HardVocabSize\TriggerMass>1.
\]  Hence the \(\HardBankSize\) candidates occupy
the first \(\HardBankSize\) marginal ranks.  For
\(k\le\HardBankSize\),
\[
 \HardMarginal(\HardCandidate{k})=\TriggerMass
 \le\HardBankSize^{-\RankTailExponent}
 \le k^{-\RankTailExponent};
\]
for \(k>\HardBankSize\), the vocabulary-floor term
\(\HardVocabSize^{-1}\) suffices.  This proves the rank--frequency inequality
in (RF\textcolor{red}{; \cref{app:eq:RF}}) with \(\FrequencyConstant=1\).  The separate vocabulary-growth clause
in (RF\textcolor{red}{; \cref{app:ass:rate-regime}}) holds whenever the public sequence \(\HardVocabSize\) has the
declared asymptotic scaling.


\emph{4. Uniform oracle accuracy (A2).}
The only nonzero oracle error is bounded by
\eqref{app:eq:hard-local-oracle-error}.  The final condition in
\eqref{app:eq:finite-hard-inclusion} therefore places it below
\((\OracleRadius)^2/(2N)\), proving (A2).  All required finite conditions hold
for every hidden matching and trigger vector.  Together with the stated
vocabulary-growth condition, this proves the asserted inclusion in the
shared asymptotic class.
\end{proof}
